\section{Eleven positive unknowns: absorb the index quotient}
\label{sec:eleven}

The second new mechanism removes an actual supplied variable from the direct
Pell construction. It gives a single equation in eleven positive integers,
of exact degree $1154$ and cost $123$ binary operations, or two equations of
maximum degree $1152$ and cost $118$ after the inherited transport shift.
The predecessor direct family uses twelve positive integers and three
equations. This is a different tradeoff from the nine-variable Jones
construction: its advantage is a moderate degree and a small explicit circuit.
The comparison is with the pinned repository collection~\cite{repo}, not a
global minimality assertion.

The old index quotient $h$ is reconstructed from a divisibility condition
inside the already quantified carrier. The first divisibility recovered is
weaker than the original one. A short-range argument for the Pell index makes
it sufficient. This is why simply deleting the old index equation would not
have worked.

\subsection{The polynomial and its exact supplied domain}

Use the entire fixed coefficient tuple $(\beta,K_0,M_C,M_F)$ of the compiler,
with $M_F$ the shifted high mask. The eleven supplied unknowns are exactly
\[
 J,F,\alpha,z,i,T,s,w,\eta,\zeta,t\in\Z_{>0}.
\]
All other letters below are polynomial abbreviations:
\begin{align*}
 q&=\beta J+1,&s_0&=s+3q,&C_w&=q-F-\alpha,\\
 X&=wq,&Y&=s_0q^3,&E&=XY,\\
 k&=\eta+\zeta,&c&=kY+\eta,&a&=E+Y,\\
 H&=4a+3,&\Delta&=a^2+H,&\ell&=EkY,\\
 A&=1+\ell(\ell+k),&G&=1+\Delta c^2,&b_0&=X+ac,\\
 I&=ic^2,&D&=1+\Delta I^2,&Q&=\Delta(D-1)=(\Delta I)^2,\\
 d&=Q-1,&v&=cT,&b&=c+RD.
\end{align*}
Here $R$ and the transport expression are
\begin{align*}
 R&=((q-1)(q-F)+\alpha)(q^2-1)+(M_C+qM_F)J,\\
 N_t&=(K_0+w)C_w+q-F-z(q-1).
\end{align*}
Define the sharp weight and new carrier by
\begin{align*}
 e&=Q(Dv^2+b^2)-1,&B_a&=2Qvb,\\
 C_*&=E+2(X+Y)+s_0q^2,&M&=I(2T+RC_*ic),\\
 L_0&=1+HM,&S&=HE,\\
 L&=EL_0,&W&=Lb_0+H(Et+k-R-1).
\end{align*}
The finite conjugate product
\begin{equation}
 \mathcal N=
 \prod_{\substack{g=\pm\sqrt G,\ \tau=\pm\sqrt A\\f=\pm\sqrt D}}
 \left[d(W-Lg-SM\tau)^2-S^2e+S^2B_af\right]
 \label{eq:eleven-norm}
\end{equation}
is an integer polynomial. Its eight formal factors define a polynomial
identity; the displayed radicals are not additional unknowns, and independent
formal signs are not assumed to be independent field automorphisms.
The single equation and two-equation form are respectively
\begin{equation}
 \boxed{N_t(1+\mathcal N)-1=0},
 \qquad
 \boxed{N_t=1,\quad\mathcal N=0}.
 \label{eq:eleven-final}
\end{equation}

For a literal division-free evaluation, put
\begin{align*}
 n&=(SM)^2A,&d_n&=dn,\\
 P_H&=S^2QDv^2,&U_H&=S^2Qb^2,\\
 e_H&=P_H+U_H-S^2,&j&=4P_HU_H,\\
 \mathcal T&=2d(W^2+L^2G),&\kappa&=[d(W^2-L^2G)]^2,\\
 \mathcal A&=2(d_n+e_H),&\mathcal B&=(d_n-e_H)^2-j,\\
 a_r&=\mathcal T-\mathcal A,&b_r&=\mathcal B-\kappa,\\
 v_r&=16d_nj-a_r(a_r\mathcal T+2b_r),&u_r&=b_r^2-a_r^2\kappa,\\
 \mathcal N&=\kappa v_r^2-u_r(\mathcal T v_r-u_r).
\end{align*}
The system compares the two final terms directly. The identity follows by
reducing a quartic modulo the quadratic with root sum $\mathcal T$ and root
product $\kappa$, exactly as in the predecessor norm; it remains valid when
quadratic fields coincide or a radicand specializes to a square.

\subsection{Recover roots without assuming the deleted equation}

At a root of the product equation, $N_t$ and $1+\mathcal N$ are integer
units. For $x_0=d(W^2+L^2G)$, one has
\[
 \mathcal T=2x_0,\qquad
 \kappa=x_0^2-4d^2W^2L^2G,\qquad
 \mathcal N\equiv(u_r-x_0v_r)^2\pmod4.
\]
Thus $1+\mathcal N\ne-1$, and $\mathcal N=0$, $N_t=1$. The authentic
transport interface now gives
\begin{equation}
 q\ge16,\quad C_w>0,\quad 3q<R<q^4-2,\quad
 Y>3q^4>3R,\quad X\ge16,\quad E>R+1.
 \label{eq:eleven-outer}
\end{equation}
Also $k\ge2$ and $kY<c<k(Y+1)$. For example, $C_w\le0$ would imply
$N_t\le1-F-(z-1)(q-1)\le0$. We do \emph{not} assume $k>E$ or $R<k$:
those predecessor bounds depended on the omitted equation.

\begin{lemma}[Uniform auxiliary bound]
Every coherent complex auxiliary branch satisfying
\[
 f^2=D,\qquad dy^2=e-B_af=Q(vf-b)^2-1
\]
obeys $|y|<M\sqrt A$ under \eqref{eq:eleven-outer}.
\end{lemma}
\begin{proof}
Since $\Delta>R+2$ and $I\ge c^2$,
\[
 Q-b=\Delta I^2(\Delta-R)-c-R>0,
 \qquad \sqrt D<\Delta I<d.
\]
For $Z=v\sqrt D+b>1$,
\[
 |y|^2\le\frac{QZ^2+1}{d}
 <[Z(1+1/d)]^2,
\]
the last numerator difference being $QZ^2-d>0$. Consequently
\begin{equation}
 |y|<v\sqrt D+b+v+1
 <cT((a+2)I+1)+c+RD+1.
 \label{eq:eleven-ybound}
\end{equation}
Set $\rho=(1+1/Y)/E$. Then $c/\ell<\rho$ and $(a+3)\rho<2$.
The exact identity
\[
 (a+2)^2\rho=E+2Y+X+6+\frac YX+
 \frac5X+\frac4Y+\frac8E+\frac4{EY}
\]
has its last four terms less than $1$. As $Y/X=s_0q^2/w\le s_0q^2$
and $X\ge16$, it is less than $C_*-2$. Divide
\eqref{eq:eleven-ybound} by $\sqrt A>\ell$, and use $(c+1)/\ell<2$:
\[
 \frac{|y|}{\sqrt A}
 <2TI+R(C_*-2)i^2c^3+2\le M.
\]
No lower bound on $k$ beyond $k\ge2$ was needed.
\end{proof}

The same outer bounds give $G>A$ and $b_0>\sqrt A$: use
$\Delta c^2>k^2Y^2(E+Y)^2>A-1$ and
$b_0=\ell+E\eta+Yc+X>\ell+k>\sqrt A$.
Every root of the inner polynomial has the form $x=M\tau+y$, hence
$|x|<2M\sqrt A$.
If a vanishing factor in \eqref{eq:eleven-norm} had $g=-\sqrt G$, then
\[
 x=\frac{W+L\sqrt G}{S}
   =t+\frac{k-R-1}{E}
     +(M+1/H)(b_0+\sqrt G)>2M\sqrt A,
\]
because the first two terms have positive sum by $E>R+1$.
This is impossible. If $G$ were nonsquare, actual quadratic conjugation
would force this forbidden branch along with the positive branch. Thus
$g$ is a positive integer and $x=(W-Lg)/S$ is rational.

In a normal splitting field, an automorphism negating an irrational $\tau$
would give $2M\sqrt A\le|y|+|y'|<2M\sqrt A$. Thus $\tau$ is rational,
hence integral. Then $y=x-M\tau$ is rational, and
$f=(e-dy^2)/B_a$ is rational, hence integral. Finally the inherited rational
Pell saturation lemma gives $y\in\Z$: if $m\ge2$, $V\in\Z$, and
$m^2V^2-(m^2-1)y^2=1$ with $y\in\mathbb Q$, then $y\in\Z$.
It follows by clearing a reduced denominator, classifying the two resulting
integer Pell solutions, and using abscissa divisibility to force ordinate
divisibility. Here apply it with $m=\Delta I$ and $V=vf-b$.
The complete argument, including zero cases, is reproduced in the accompanying
proof file. In particular $x$ is integral. The carrier equality implies
\begin{equation}
 E\mid H(k-R-1).
 \label{eq:eleven-weakdiv}
\end{equation}
At this stage it would be invalid to cancel $H$.

\subsection{Recover the main index before the index quotient}

Let $A_0=a+2$, and write $\chi_{A_0},\psi_{A_0}$ for the usual Pell
sequences. The fundamental positive integer-coordinate solution for
$A_0^2-1$ is $(A_0,1)$. Main and normalized strong norm recovery gives
\[
 g=\chi_{A_0}(p),\quad c=\psi_{A_0}(p),\quad
 |f|=\chi_{A_0}(m),\quad I=\psi_{A_0}(m),\qquad p,m\ge1.
\]
The standard strong-divisibility and composition identities now give
\begin{equation}
 p\mid m,\qquad c\mid m,\qquad m\ge pc>2p.
 \label{eq:eleven-rank}
\end{equation}
Here is the direct argument, to make its independence from the deleted
equation explicit. Since $c^2\mid I$, strong divisibility first gives
$p\mid m$. Write $m=pj$. Then
$\psi_{A_0}(pj)=c\psi_{\chi_{A_0}(p)}(j)$, while modulo $c$ the
second factor is $j\chi_{A_0}(p)^{j-1}$. Since
$\chi_{A_0}(p)^2\equiv1\pmod c$, we get $c\mid j$. No ordinary-strong
large-$c$ rank lemma is necessary. Also $c>2Y>6R$, and $p\ge2$ with
$c\ge(2A_0-1)^{p-1}>2p$.

The inherited signed auxiliary-index argument can therefore be applied
before the main congruence modulo $H$. For clarity, its essential steps are
included. Put $S_0=\Delta I$ and $V=c(Tf-1)-Rf^2$. The auxiliary Pell
equation shows that
\[
 S_0|V|=\chi_{S_0}(j_0),\qquad |y|=\psi_{S_0}(j_0)
\]
for a positive odd $j_0=2r+1$. The usual odd-index polynomial $Q_r$ has
\[
 \chi_{S_0}(2r+1)=S_0Q_r(S_0^2),\quad
 Q_r(1-A_0^2)=(-1)^r\psi_{A_0}(2r+1),\quad
 Q_r(0)=(-1)^r(2r+1).
\]
Because $Q=\Delta(f^2-1)\equiv1-A_0^2\pmod{|f|}$ and
$V\equiv-c\pmod{|f|}$, these give
$\psi_{A_0}(j_0)\equiv\pm c\pmod{|f|}$. Squaring and doubling yields
\[
 \chi_{A_0}(2j_0)\equiv\chi_{A_0}(2p)\pmod{\chi_{A_0}(m)}.
\]
Signed index reduction, using $0<2p<m$, gives $j_0\equiv\pm p\pmod m$,
hence modulo $c$. This follows by reducing modulo $2m$, reflecting indices
to $[0,m]$, and using strict growth and
$2\chi_{A_0}(m-1)<\chi_{A_0}(m)$ to exclude a negative representative.
Meanwhile $Q\equiv0\pmod{c^4}$ and $V\equiv-R\pmod c$ give
$j_0\equiv\pm R\pmod c$. Since $p,R\in(0,c/2)$, one obtains
\begin{equation}
 p=R.
 \label{eq:eleven-pR}
\end{equation}

\subsection{A short index range restores the full divisibility}

The first norm has fundamental unit $(2EY+1,2)$, so for $P_0=2EY+1$,
\[
 |\tau|=\chi_{P_0}(n),\qquad k=2\psi_{P_0}(n),\qquad n\ge1.
\]
Because $P_0>A_0$ and $c>k$, we have $n<p=R$. Put
$g_0=\gcd(E,H)$ and $\mu=E/g_0$. Since $H=4XY+4Y+3$,
\begin{equation}
 \gcd(Y,H)=\gcd(Y,3)\le3,\qquad
 g_0\mid X\gcd(Y,H),\qquad \mu\ge Y/3>R.
 \label{eq:eleven-modulus}
\end{equation}
Equation \eqref{eq:eleven-weakdiv} gives $k\equiv R+1\pmod\mu$.
As $P_0\equiv1\pmod E$, the Pell recurrence also gives
$k\equiv2n\pmod E$. Hence $2n\equiv R+1\pmod\mu$. But
$1\le n\le R-1$ gives $|2n-R-1|<R<\mu$. Therefore
\[
 2n=R+1,\qquad E\mid(k-R-1).
\]
Set $h=(k-R-1)/E$. It is integral and positive: $R>3q\ge48$ gives
$n\ge25$ and $k=2\psi_{P_0}(n)>2n=R+1$.

Divide the carrier identity by $E$ to get
\[
 L_0b_0+H(t+h)=L_0g+H(M\tau+y).
\]
Since $L_0\equiv1\pmod H$, we obtain $g\equiv b_0\pmod H$.
Thus $\gamma=(g-b_0)/H$ is an integer. Moreover $n\ge2$ gives
$k\ge4P_0>X$, hence $c>X$, and
$g-b_0>(a+1)c-ac-X=c-X>0$. So $\gamma>0$.
All original normalized compiler equations are now restored, including
$k-hE-R=1$ and the main congruence. The signed-quotient soundness theorem
handles a possible negative $f$ by changing $f,T$ together, and extracts a
blank K215 halt. The order of recovery above avoids every use of $h$ before
its reconstruction.

\subsection{Completeness and exact degree}

For a canonical positive compiler witness the inherited theory gives
$X=2^R$, $R>3q$, and $Y>2^R$. Since $q\ge16$ and
$2^{3q}>3q^4$, its old coordinate $s_{\rm old}=Y/q^3$ satisfies
$s_{\rm old}>3q$. Hence $s=s_{\rm old}-3q>0$ is legitimate. Define
\[
 t=L_0\gamma+M\tau+y-h.
\]
This is positive because $h<k/E$ and $\tau>\ell=EkY>k/E$.
It makes a coherent factor in \eqref{eq:eleven-norm} vanish and preserves
$N_t=1$. This proves both directions of existence equivalence.
The shifted system uses
\[
 (K_0+w)C_w=F+z(q-1),
\]
with $z_{\rm old}=z+1$. Canonical witnesses have $w=2^R/q>q>F$ and
$C_w>0$, so $z_{\rm old}\ge2$. This saves one gate while preserving
positivity and existence.

For the degree calculation, let $\mathcal N_{\rm old}$ be the predecessor's
sharp half-carrier norm, of degree $1056$. In the localization at $E$,
\begin{equation}
 \mathcal N=E^{16}\mathcal N_{\rm old}
       \left(t+\frac{k-R-1}{E}\right),\qquad s_{\rm old}=s+3q.
 \label{eq:eleven-degreeidentity}
\end{equation}
Each of the eight factors gains $E^2$, and its degree in the old carrier
is at most $16$, so the denominators clear. The rational shift has total
degree $-2$, below $\deg t=1$, and cannot change a highest homogeneous part.
Writing $q_0=\beta J$, $k_0=\eta+\zeta$, and $s_{00}=s+3q_0$, the
predecessor's proved leading form gives
\[
 \mathcal N_{1152}
 =2^{46}i^{48}k_0^{96}w^{90}s_{00}^{192}
 q_0^{710}(q_0-F)^{16}\ne0.
\]
Thus $\deg\mathcal N=1056+16\cdot6=1152$. The transport leading form is
$w(q_0-F-\alpha)-zq_0\ne0$, proving exact product degree $1154$ and
system residual degrees $2,1152$, uniformly for the authentic $\beta>0$.

The delivered circuit has $75$ multiplications and $48$ additions or
subtractions for the product. Independent checks include $90$ formal norm
products modulo three primes, $77$ carrier-denominator identities, exact
integer and rational fixtures, $500$ parity fixtures, and two complete
univariate polynomial calculations of degree and leading coefficient.
Another $320$ bound fixtures deliberately satisfy $k\le E$, testing the
hypothesis removed from the predecessor argument. The four actual fixed
payloads are hashed and used in $40$ additional modular norm checks.
These diagnostics support the transcription; the unbounded proof above and
the inherited compiler calibration establish the theorem in ZFC.
